\chapter{The Nature of Data}
\label{ch:nature-of-data}

\begin{chapterintro}
Before a single line of code runs, before a CPU executes a single instruction,
before a bit flips in memory, there is data. This chapter asks a question
most programmers skip: what is data, exactly? Not how to store it, not how
to query it, but what it fundamentally is. The answer matters more than it
first appears to. The way you think about data shapes every decision you
make when working with it.
\end{chapterintro}
{\small \linespread{1.5}\selectfont
We no longer suffer from the weight of history, the way Marx once put it.
Instead, we suffer from the burden of an ever-growing accumulation of
information \autocite{ratzky2005}. Given the sheer volume of data that
saturates contemporary society, an answer to almost any question is only a
few nanoseconds away, regardless of where that question is asked
\autocite{mccal2004}. Because answers are so easily available, it is
understandable that many members of what is loosely called the
\term{millennial generation} have begun to question the value of a research
paper, or the point of an assignment, at all \autocite{tansen2013}.\\
That question deserves a real answer, and a real answer has to start
further back than it first appears to. Before any discussion of types of
data can be useful, a more basic confusion needs clearing up: everyday
usage stretches the word ``data'' loosely enough that it ends up meaning
several different things depending on who is using it. This chapter is, in
large part, an attempt to slow down long enough to say what the word
actually commits us to, not a survey of data types, and not yet a theory of
how data are stored, only an answer to what counts as data at all.\\
The title of this chapter, ``the nature of data,'' is fairly broad. What we
are actually discussing here are the fundamentally different ways of
approaching data, and of working with it once you have it. The discussion
that follows may read as somewhat academic, and dry in places, but the view
it builds is the one the rest of this book keeps coming back to, all the
way down to the last byte and the last array.}
\newpage

% ============================================================
\section{What Is Data?}
\label{sec:what-is-data}
% ============================================================

\begin{sectionintro}
The word \textit{data} appears everywhere in computing, in textbooks, job
titles, product names, yet it is rarely defined with any precision. Before
we can talk about how computers represent and process data, we need a
working definition that is precise enough to be useful.
\end{sectionintro}
A simple example is the best place to start. Looking at the trunk of a tree
reveals facts about it without anyone first explaining what a tree trunk
is. The number of growth rings gives a reasonable estimate of the tree's
age, since each ring corresponds to one year of growth, laid down during
the annual cycle of dormancy and activity. The texture and color of the
bark suggest its species; warped grain, discoloration, or insect tunnels
point toward its health. Look more closely still, and the wood itself, its
density, the direction of its grain, its moisture content, accounts for how
much that piece of wood weighs and how much space it occupies. None of
these facts were created the moment they were noticed; they existed in the
tree itself, available to anyone who knew how to read them. That
correspondence, between a physical state and a fact about the world that
can be read off from it, is the simplest possible illustration of what
this chapter means by \term{data}.\\
The word \term{data} comes from the Latin \textit{datum}, the past
participle of \textit{dare}, ``to give.'' A datum is, quite literally, a
``given'': a piece of information handed over, taken as a starting point
for reasoning rather than derived from something else. In everyday speech,
the plural form ``data'' is often treated as a singular collective noun
(``the data is clear''), but in technical writing the plural is preferred:
data \textit{are} observations, measurements, or recorded facts, not a
single undivided thing. For the sake of consistency, this book treats
\textit{data} as a plural throughout.
\vfill
\begin{historical}
The Latin origin is not merely a piece of trivia. It carries real weight
that survives into technical usage: a datum is something received from
the world, not invented by the observer who records it.
\end{historical}
\newpage
{\small \linespread{1.2}\selectfont
Different sources define the term in different, only partially overlapping,
ways. Rob, Morris, and Coronel describe data as \term{raw facts}: facts
that have not yet been processed for their meaning to become apparent
\autocite{rob2013}. Ratzan defines the companion term instead:
\term{information} is a coherent collection of data, organized in a
particular way and meaningful as a result \autocite{ratzan2004}. The two
definitions agree on more than they disagree about: in both, data precede
meaning rather than supply it directly.\\
General dictionaries describe much the same shape from a different angle.
Webster's \textit{New International Dictionary} treats data as something
given or accepted: facts or principles granted, the basis from which an
inference or a line of reasoning proceeds, or the material out of which an
idealized system of any kind is built \autocite{websters1934}. The
\textit{Oxford English Dictionary} is more compact: data are facts or
things known, used as the basis for inference or calculation
\autocite{oed2023}. Both dictionaries note, in passing, the detail
mentioned above: ``data'' is grammatically the plural of ``datum,'' yet it
is used so often as a singular, collective noun that treating it as
singular is equally acceptable in ordinary writing.

\begin{definition}[Data, UNESCO]
Facts, concepts, or instructions expressed in a standardized way, suitable
for being communicated, interpreted, or processed, whether by people or by
automated systems \autocite{unesco2008}.
\end{definition}
Taken together, these definitions point at something none of them states
outright. Data have no shape that belongs to them independently of some
particular point of view. A fact only counts as data if it sits within a
framework, a framework of meaning, of purpose, or of use, relative to
which it is relevant, organized, coherent, and useful. A growth ring means
nothing to a system with no concept of seasons; a voltage means nothing to
a circuit built to ignore it. Data presuppose some prior agreement, however
informal, about what is worth attending to in the first place.

\begin{example}[The Same Bytes, Different Meanings]
The five bytes \texttt{0x41 0x52 0x4C 0x49 0x5A} can be interpreted in several
different ways:

\begin{itemize}
    \item the ASCII string \texttt{"ARLIZ"}
    \item five separate unsigned 8-bit integers:
          65, 82, 76, 73, and 90
    \item the first five bytes of some binary file format
    \item a sequence of bytes transmitted over a network connection
\end{itemize}

The bytes themselves are identical in every case. What changes is not the
underlying data, but the interpretation imposed by context. Meaning is not
contained in the bytes themselves; it arises from the agreement between the
system that produced them and the system that reads them.
\end{example}
}
\newpage
% ------------------------------------------------------------
\subsection{Data in Mathematics}
\label{subsec:data-in-mathematics}
% ------------------------------------------------------------

\begin{subsectionintro}
The definitions above are useful, but they are also informal, the same
word doing slightly different work in each one. Mathematics offers a
sharper handle on the same idea, and it is the handle this book keeps
returning to.
\end{subsectionintro}
In mathematics, data are most naturally thought of as a \textit{function}:
a mapping from some index set to some value set.

\begin{definition}[Datum, formal]
A \term{datum} is a value drawn from some set of possible values, called
its \term{domain}. A \term{collection of data} is most simply modeled as a
function from an \term{index set} into a \term{value set}: each index
names a position, and the function returns the value held at that
position, and only at that position.
\end{definition}
This framing looks abstract, but it matches ordinary cases closely. A
single temperature reading is a datum drawn from the real numbers; a
calendar date is a datum drawn from a finite set of valid dates; a pixel's
color is a datum drawn from the set of representable RGB triples. A weather
station's log of one reading per hour is, formally, a function from
\textit{hour number} (the index set) to \textit{temperature} (the value
set), which is already, without yet using the word, an array. That is not
a coincidence: a closely related definition, a function from an index set
to a value set, is standard throughout the mathematical literature on
sequences and tuples \autocite{knuth1997}, and starting with the next
chapter this book uses exactly that definition as the formal starting
point for arrays themselves.\\
This framing is more than abstract elegance. It tells you immediately what
operations make sense on a piece of data. A function from indices to
values supports indexing (``give me the value at position 3'') and
iteration (``apply this operation to every value''). It does not, by
itself, support operations that require knowing the \textit{type} of the
values; for that, you need to know more about the value set itself. The
mathematical view also makes precise what it means for two pieces of data
to be the \textit{same}: two arrays are equal only if they have the same
index set and the same value at every index. Equality, in this view, is
not about appearance; it is about structure. The mathematical view does
not replace the informal ``raw facts'' definition; it gives that
definition a shape precise enough to build on.
\newpage
% ------------------------------------------------------------
\subsection{The Perishability of Data}
\label{subsec:perishability-of-data}
% ------------------------------------------------------------

\begin{subsectionintro}
Data are not a permanent, objective record of the world. They are a
snapshot taken at a particular moment, by a particular observer, using a
particular method. Understanding this keeps you from over-trusting your
data, and from under-appreciating the engineering required to make it
reliable.
\end{subsectionintro}
A common assumption, especially early in one's programming career, is that
data are fixed and objective: you record a fact, and the fact is stored.
This is only partially true, and the part that is not true matters.
Walliman observes that data are not only imprecise but \term{transient}
and \term{perishable} \autocite{walliman2011}. A poll measuring voting
intentions ahead of an election will not produce the same result as an
otherwise identical poll conducted at a different time, even with the same
sampling method and, in principle, the same respondents
\autocite{pippenborg2011}.
\begin{notebox}
``Perishable'' here does not mean that data physically vanish. A
measurement recorded yesterday is still there to be read today. What
perishes is its \textit{currency}: its claim to describe the present state
of whatever it measured.
\end{notebox}
Data perish in a second sense as well. Secondhand reports and biased
accounts are routinely presented as settled fact, because no method of
recording information is perfectly reliable, and some distortion of
interpretation enters at almost every stage of transmission. This is not
a rare edge case confined to opinion polling; it is a fundamental property
of how facts move from the world into a recorded form, and a significant
portion of this book is dedicated to the engineering that addresses it
(error detection, error correction, checksums, and so on).
\begin{notebox}
None of this means data are unreliable or useless. It means that
reliable, useful data require careful engineering: choosing the right
precision, designing for integrity in transmission, recording provenance
and timestamp, and being explicit about what was measured and how.
\end{notebox}
On the question of how data connect to the rest of what we know, Walliman
describes data as related to knowledge through a hierarchical system: data
sit at the general, abstract end of that hierarchy, and become more
particular and more concrete only as they are organized, interpreted, and
put to use \autocite{walliman2011}. That hierarchy is the subject of the
next section.
\newpage
% ============================================================
\section{Data, Information, and Knowledge}
\label{sec:data-information-knowledge}
% ============================================================
{\small \linespread{1.1}\selectfont
\begin{sectionintro}
The words \textit{data}, \textit{information}, and \textit{knowledge} are
used interchangeably in casual conversation, but they describe three
distinct things. The distinction matters in practice: confusing them
leads to building systems that produce numbers when people need answers.
\end{sectionintro}
The relationship between data, information, and knowledge is often
described as a hierarchy. We will not attempt a complete account of that
hierarchy here, but the first two transitions are worth understanding
precisely, since they are the ones this book depends on.\\
Return to the tree. A raw count of growth rings is data: fifty-three
concentric lines, nothing more. Data are syntactic; they have structure,
but not yet meaning. Noticing that fifty-three rings means the tree is
roughly fifty-three years old, and that the unusually narrow rings around
year thirty correspond to a known regional drought, turns that count into
\term{information}: data organized and placed in context, in keeping with
Ratzan's definition above \autocite{ratzan2004}. Information answers
questions like ``what happened?'' and ``how much?''. It is still factual
and objective, but it is now interpretable by a human being. Knowing how
to read ring width as a record of annual rainfall, reliably enough to
apply the same method to a tree no one has studied before, is
\term{knowledge}: information understood thoroughly enough to be put to
new use.

\begin{example}[From Data to Knowledge]
\begin{itemize}
    \item \textbf{Data:} \texttt{39.5}
    \item \textbf{Information:} Patient temperature is 39.5\textdegree C, recorded
          at 14:32.
    \item \textbf{Knowledge:} This patient has a significant fever. Combined with
          the readings from the past six hours, the temperature is rising. The
          attending physician should be notified.
\end{itemize}
\end{example}
The three terms are often used as if they were interchangeable, and the
confusion this causes is part of what motivates the observation, mentioned
earlier, that easy access to information makes the work of building
knowledge look optional \autocite{tansen2013}. It is not. A search engine
can return information in a few nanoseconds \autocite{mccal2004}; it
cannot, by itself, supply the understanding required to know whether that
information is relevant, reliable, or correctly applied. The gap between
having information and having knowledge is exactly the gap that instant
retrieval does not close.

\begin{notebox}
This is a simplification of a genuinely contested area. Some writers add
a fourth tier, wisdom, above knowledge, and some order the lower tiers
differently. That debate is not settled here, and it does not need to be:
the three-term version above is enough to keep data, information, and
knowledge from collapsing into one another for the rest of this book.
\end{notebox}
}Why does this matter for a book about arrays and data representation?
Because computers operate almost entirely at the level of \textit{data}.
A CPU does not understand information; it manipulates bits. A memory array
does not know what its contents mean; it stores and retrieves bytes at
addresses. The transition from data to information is something that
software engineers design into their systems deliberately, through type
systems, schemas, protocols, and documentation. When you write
\texttt{int temperature = 39;}, you are asserting that this particular
sequence of bits represents a temperature, measured in some unit, with
some precision. The meaning is in your head and in your documentation, not
in the machine. The machine stores \texttt{00000027} and nothing else.\\
This is not a criticism of computers. It is a description of how they
work, and understanding it is the first step toward understanding why the
rest of Volume~I, all the encoding schemes, type systems, and
representation formats, exists. Each one is an answer to the question: how
do we encode information, which has meaning, as data, which does not, in a
way that preserves the meaning reliably?
\newpage
% ============================================================
\section{The Relationship Between Data and Representation}
\label{sec:data-and-representation}
% ============================================================

\begin{sectionintro}
Data and representation are two different things, easy to blur together.
This section draws the line between them; the next chapter takes
representation up as a subject in its own right.
\end{sectionintro}
Consider the number seven. The number seven is an abstract mathematical
object; it exists independently of how we write it, speak it, or store
it. To communicate or store it, we must choose a \textit{representation}:
the digit \texttt{7}, the Roman numeral \texttt{VII}, the binary string
\texttt{0111}, three fingers held up. Each of these stands for the same
abstract value; none of them \textit{is} that value.\\
The same separation holds for data. A growth ring is a fact about a tree
whether or not anyone counts it; the fact only becomes data, in the sense
this chapter has been building toward, once it is expressed in some form
that can be communicated or processed, Roman numerals, binary, notches on
a tally stick, a single byte in a file. Data are the facts; a
representation is whatever has been chosen to carry them. The choice is
never forced by the data alone, and changing the representation never
changes what the data are about.\\
That is as far as this chapter needs to go. \textit{What} a representation
is, formally, and what makes one representation better than another, is
the subject of Chapter~\ref{ch:philosophy-of-representation}.
\bigskip
\noindent\textit{The next chapter asks a more fundamental question: what
does it mean to represent something at all? That question turns out to
have a surprisingly deep answer.}